% =============================================================================
%  senior_project_description.tex
%  Undergraduate Senior Project — one-paragraph description for UnionLabs.
%
%  Compile:  pdflatex senior_project_description.tex
%  (Uses only article + geometry, so it builds on a bare TeX install and on
%   Overleaf. No parskip / enumitem / hyperref.)
%
%  Lives in its OWN folder, not directly in docs/, because docs/*.tex is
%  gitignored there (pandoc writes LaTeX intermediates next to the PDFs it
%  builds). This is a source document, not a build artifact.
%
%  The paragraph follows four beats:
%      (1) what UnionLabs is
%      (2) what is left to design   -- web / hardware integration / experiments
%      (3) the undergraduate mission
%      (4) what the students do first
% =============================================================================
\documentclass[11pt]{article}
\usepackage[margin=1in]{geometry}
\setlength{\parindent}{0pt}

\begin{document}

\begin{center}
{\large\bfseries UnionLabs: AWS-based Remote Access and Sharing of NextG
and IoT Testbeds}
\end{center}

\vspace{1em}

UnionLabs is a cloud-based platform that makes wireless research testbeds
remotely accessible and shareable, so that experiments on NextG and IoT
hardware can be run from a browser by users at any institution. The core
mechanism is already prototyped: a reservation is translated automatically
into an isolated environment with the reserved radios attached, and is served
to the user as a remote desktop. Three components remain to be built on top of
it: a web front end providing user accounts, a catalog of available resources,
a reservation calendar, and session management; hardware integration that
brings additional device classes, in particular IoT nodes, under the same
reservation model; and a library of ready-to-run wireless experiments that a
remote user can launch to verify a reservation. Most of the platform is
already in place, so a senior project team would build on existing work rather
than start from scratch. The students would first study the existing system
and codebase, covering Python, JavaScript, and C++, together with the basic
wireless background. They would then take the web design as an independent
mission, with the implementation of wireless experiments as an optional
additional task.

\end{document}
